\documentclass{aastex631}

\usepackage{amsmath,amssymb}
\usepackage{natbib}

\newcommand{\dd}{\mathrm{d}}
\newcommand{\Jt}{\tilde{J}}
\newcommand{\Jtp}{\tilde{J}'}
\newcommand{\Jtpp}{\tilde{J}''}
\newcommand{\Req}{R_{\rm eq}}
\newcommand{\Rpol}{R_{\rm pol}}

\shorttitle{ORCHARD CMS Gravity Methods}
\shortauthors{Tejada Arevalo}

\begin{document}

\title{The Concentric Maclaurin Spheroid Gravity Solver in ORCHARD:
Physics, Methodology, and Application to the Gravitational Observables of
Jupiter and Saturn}

\author{Roberto Tejada Arevalo}

\begin{abstract}
We describe the Concentric Maclaurin Spheroid (CMS) gravity solver
implemented in the ORCHARD planetary evolution code
(\texttt{utils/cms\_hubbard.py}).  The CMS method of \citet{Hubbard2013}
represents a rotating fluid planet as $N$ nested spheroids of constant
density whose surfaces are iterated to equipotentials of the combined
gravitational and centrifugal potential.  Unlike the perturbative Theory
of Figures (ToF), which expands the figure functions in powers of the
smallness parameter of rotation, CMS is nonperturbative: its only error
sources are the radial discretization of the density profile and the
angular quadrature, both of which are controllable.  We present the
multipole expansion of the spheroid-stack potential, the equipotential
shape equations, and the quadrature formulae for the zonal harmonics
$J_{2n}$, the equatorial and polar radii, the oblateness, and the axial
moment of inertia, in the conventions of \citet{Hubbard2013} and
\citet{MilitzerHubbard2024} as implemented.  The ORCHARD implementation
is a clean-room realization of the published equations, accelerated with
a vectorized Newton shape solve, layer-summed coefficient recurrences,
\texttt{numba} JIT kernels, and warm-started Jacobi iteration, which
together bring the per-solve cost to within a factor of
$\sim$1.5--1.9 of ToF7.  We describe the mass-conserving mean-to-equatorial
radius mapping that lets CMS drop into ORCHARD's existing gravity-backend
dispatcher (\texttt{gravity\_method = cms}), the validation against the
analytic Maclaurin spheroid and the published Jupiter models of
\citet{MilitzerHubbard2024}, and the two ways CMS is used to compute the
gravitational observables of Jupiter and Saturn: as an in-loop backend
during evolution, and as the ``CMS shell-consistent'' post-processing
score applied to evolution grids, including the spin self-consistency
loop and the Saturn gravity reference-radius rescale.
\end{abstract}

\keywords{planets and satellites: interiors --- planets and satellites:
gaseous planets --- methods: numerical --- gravitation}

% ======================================================================
\section{Introduction}\label{sec:intro}
% ======================================================================

The zonal gravitational harmonics $J_{2n}$ of Jupiter and Saturn,
measured to extraordinary precision by Juno \citep{Iess2018} and by the
Cassini Grand Finale \citep{Iess2019}, are the primary observational
constraints on giant-planet interior structure.  A rotating fluid planet
in hydrostatic equilibrium acquires an oblate figure, and its external
potential departs from spherical symmetry:
\begin{equation}
U_{\rm grav}(r,\mu) \;=\; \frac{GM}{r}\left[\,1 -
\sum_{n=1}^{\infty} \left(\frac{a}{r}\right)^{2n} J_{2n}\,
P_{2n}(\mu)\right],
\label{eq:external}
\end{equation}
where $a$ is the equatorial reference radius, $\mu = \cos\theta$ with
$\theta$ the colatitude, and $P_{2n}$ are Legendre polynomials.  The
harmonics are weighted integrals of the density distribution over the
rotationally distorted figure; computing them from a model interior
requires solving for the figure and the potential self-consistently.

Two families of methods dominate.  The perturbative
\emph{Theory of Figures} \citep[ToF;][]{ZharkovTrubitsyn1978} expands the
level surfaces in powers of the rotational smallness parameter
$m = \omega^2 s^3 / GM$ (with $s$ the mean radius); ORCHARD's default
gravity backends are the fourth-order scheme of \citet{Nettelmann2017}
(\texttt{utils/TOF.py}) and the seventh-order scheme of
\citet{Nettelmann2021} (\texttt{utils/TOF7.py}).  The nonperturbative
\emph{Concentric Maclaurin Spheroid} method
\citep[CMS;][]{Hubbard2013} instead discretizes the planet into nested
constant-density spheroids and iterates their shapes to exact
equipotentials, with no expansion in the rotation rate.  CMS is the
method behind the Juno-era Jupiter models of \citet{Wahl2017} and
\citet{MilitzerHubbard2024} and the Saturn models of \citet{Militzer2019},
so implementing it in ORCHARD also buys direct methodological
comparability with those works.  The trade-offs between the two families
were quantified for ORCHARD in the \citet{MilitzerHubbard2024}
reproducibility study (\texttt{validation/mh24/REPORT.md}): on identical
density profiles, ToF7 and CMS agree on $J_2$ to
$\sim\!10^{-5}$ (coreless) to $\sim\!10^{-4}$ (compact core) relative ---
far below current EOS uncertainties, but \emph{above} Juno's relative
precision on $J_2$ of $\sim\!4\times10^{-8}$.  For $J_2$-precision work
and for apples-to-apples comparison with the CMS literature, the
nonperturbative solver is therefore worth having.

This document describes the physics and the numerics of ORCHARD's CMS
implementation, \texttt{utils/cms\_hubbard.py}.  It is a clean-room
implementation written from the published equations of
\citet{Hubbard2013} (his Eqs.~40--52) and \citet{MilitzerHubbard2024}
(their Eqs.~9--13), \emph{not} a port of the GPL-licensed
\texttt{nmovshov/CMS-planet} code (which was consulted only to
cross-check equation forms); this keeps the module licensable with the
rest of ORCHARD.  Section~\ref{sec:physics} sets up the physical picture
and conventions.  Sections~\ref{sec:potential}--\ref{sec:observables}
give the multipole expansion, the equipotential shape equations, and the
formulae for the harmonics and shape observables.
Section~\ref{sec:numerics} details the numerical scheme and its
accelerations, Section~\ref{sec:integration} the integration into
ORCHARD's evolution loop, Section~\ref{sec:validation} the validation
suite, and Section~\ref{sec:observables_js} the application to the
observed gravitational moments of Jupiter and Saturn, including the
post-processing scoring pipeline.  Section~\ref{sec:limitations} lists
the current limitations.

% ======================================================================
\section{Physical Picture and Conventions}\label{sec:physics}
% ======================================================================

A single homogeneous, rigidly rotating, self-gravitating fluid body in
equilibrium is a \emph{Maclaurin spheroid}: an oblate spheroid whose
eccentricity is an exact, classical function of its rotation rate
\citep{Chandrasekhar1969}.  The CMS method generalizes this to an
arbitrary barotropic density profile by superposing $N$ \emph{concentric}
spheroids of constant density.  A continuous profile $\rho(r)$ is
represented as a stack of density increments: spheroid $i$ carries the
density \emph{jump}
\begin{equation}
\delta_i \;=\; \rho_i - \rho_{i-1}, \qquad \rho_{-1} \equiv 0,
\label{eq:delta}
\end{equation}
where $\rho_i$ is the physical density of the material between surface
$i$ and surface $i+1$.  Summing the increments interior to any point
recovers the local physical density, so the stack of uniform spheroids
reproduces the layered planet exactly.  Each surface is then required to
be an equipotential of the total (gravitational plus centrifugal)
potential; because each layer between adjacent equipotentials has
constant density, the configuration is a consistent discrete
representation of a barotrope in hydrostatic equilibrium.  The
nomenclature and conventions, following \citet{Hubbard2013} and
implemented verbatim in the module header, are:
\begin{itemize}
\item Spheroids are indexed $i = 0$ (outermost) through $N-1$
  (innermost) --- the same surface-to-center ordering as ORCHARD's mass
  grid.
\item $\lambda_i = a_i / a_0$ is the equatorial radius of surface $i$
  normalized to the outer equatorial radius, so $\lambda_0 = 1$ and the
  sequence is descending.
\item $\zeta_i(\mu) = r_i(\mu)/a_i$ is the \emph{shape function} of
  surface $i$: its radius at $\mu = \cos\theta$, normalized to its own
  equatorial radius.  By construction $\zeta_i(0) = 1$ at the equator and
  $\zeta_i < 1$ toward the pole for an oblate figure.
\item The rotation parameter is
  \begin{equation}
  q \;=\; \frac{\omega^2 a_0^3}{GM},
  \label{eq:q}
  \end{equation}
  built on the \emph{equatorial} radius (contrast ToF's
  $m = \omega^2 s^3/GM$ on the mean radius; for Jupiter $q \simeq 0.089$
  while $m \simeq 0.083$).
\item North--south symmetry is assumed, so all integrals run over
  $\mu \in [0,1]$ with even Legendre polynomials only, and only even
  harmonics $J_{2n}$ are produced.
\end{itemize}
The solution unknowns are the $N \times L$ values of $\zeta_i(\mu_a)$ on
a Gauss--Legendre grid of $L$ colatitude nodes; the harmonics follow from
quadratures over the converged shapes.

% ======================================================================
\section{Potential of the Spheroid Stack}\label{sec:potential}
% ======================================================================

The total potential at a point on surface $j$ receives three kinds of
contribution: (i) spheroids at or interior to $j$ ($i \ge j$), seen from
outside, contribute through their \emph{exterior} multipole expansions;
(ii) spheroids exterior to $j$ ($i < j$), within which the point lies,
contribute through their \emph{interior} expansions, which contain both
$r^{2k}$ terms and an $r^2$ term; and (iii) the centrifugal potential
$\tfrac{1}{2}\omega^2 r^2 (1-\mu^2)$.  Following \citet[][Eqs.~40--43]
{Hubbard2013}, the dimensionless multipole moments of each spheroid are
defined from its shape by quadratures (implemented in
\texttt{\_multipoles}):
\begin{align}
\Jt_{i,2k} &= -\frac{3}{2k+3}\,
  \frac{\delta_i \lambda_i^3}{D} \int_0^1 P_{2k}(\mu)\,
  \zeta_i^{\,2k+3}(\mu)\,\dd\mu ,
\label{eq:Jint}\\[2pt]
\Jtp_{i,2k} &= -\frac{3}{2-2k}\,
  \frac{\delta_i \lambda_i^3}{D} \int_0^1 P_{2k}(\mu)\,
  \zeta_i^{\,2-2k}(\mu)\,\dd\mu , \qquad k \ge 2,
\label{eq:Jext}\\[2pt]
\Jtp_{i,2} &= -3\,\frac{\delta_i \lambda_i^3}{D}
  \int_0^1 P_{2}(\mu)\, \ln \zeta_i(\mu)\,\dd\mu ,
\label{eq:Jext2}\\[2pt]
\Jtp_{i,0} &= -\frac{3}{2}\,\frac{\delta_i \lambda_i^3}{D}
  \int_0^1 \zeta_i^{\,2}(\mu)\,\dd\mu ,
\label{eq:Jext0}\\[2pt]
\Jtpp_{i} &= \phantom{-}\frac{1}{2}\,\frac{\delta_i \lambda_i^3}{D},
\label{eq:Jpp}
\end{align}
where the common normalization
\begin{equation}
D \;=\; \sum_{i=0}^{N-1} \delta_i \lambda_i^3 \int_0^1
\zeta_i^{\,3}(\mu)\,\dd\mu
\label{eq:D}
\end{equation}
is proportional to the total mass ($M = 4\pi a_0^3 D$), so that all
moments are mass-fraction--weighted and the overall scale of $\rho$
drops out --- only density \emph{ratios} matter.  The $k=1$ exterior
moment picks up the logarithm because the generic exponent $2-2k$
vanishes there, and $\Jtpp_i$ is the coefficient of the $r^2$ piece of
the interior potential of a uniform spheroid.  Note the leading minus
signs: for a physically positive mass increment, $\Jt_{i,0}$ is
\emph{negative}, and in particular $-\sum_{i\ge j}\Jt_{i,0}\,
(\lambda_i/\lambda_j)^0$ evaluates to the enclosed mass fraction
interior to surface $j$.

For the shape iteration it is convenient to collect, for every surface
$j$ and every order $k$, the layer-summed coefficients (implemented in
\texttt{\_coeffs}):
\begin{align}
A_{j,k} &= \sum_{i \ge j} \Jt_{i,2k}
  \left(\frac{\lambda_i}{\lambda_j}\right)^{2k}
  &&\text{(interior spheroids)},
\label{eq:A}\\
B_{j,k} &= \sum_{i < j} \Jtp_{i,2k}
  \left(\frac{\lambda_j}{\lambda_i}\right)^{2k+1}
  &&\text{(exterior spheroids)},
\label{eq:B}\\
C_{j} &= \sum_{i < j} \Jtpp_{i}
  \left(\frac{\lambda_j}{\lambda_i}\right)^{3}
  &&\text{(quadratic term)},
\label{eq:C}
\end{align}
together with the outer-surface special case
$A^{(0)}_k = \sum_i \Jt_{i,2k}\lambda_i^{2k}$ used in
Equation~(\ref{eq:outer}) below.  Because the sums over $i \ge j$ and
$i < j$ are cumulative, all $N$ surfaces' coefficients are obtained at
once by one reverse and one forward cumulative sum per order --- an
$\mathcal{O}(NK)$ operation ($K = k_{\max}+1$ retained orders) that makes
the subsequent per-surface, per-latitude equation independent of $N$.
This is the same layer-summation idea that underlies the ``accelerated''
CMS of \citet{Militzer2019}.

% ======================================================================
\section{Equipotential Shape Equations}\label{sec:equipotential}
% ======================================================================

Let $U(r,\mu) = U_{\rm grav} + \tfrac{1}{2}\omega^2 r^2 (1-\mu^2)$ be
the total potential.  Nondimensionalizing on surface $j$ by the factor
$\lambda_j a_0 / GM$ and using
$\tfrac{1}{2}(1-\mu^2) = \tfrac{1}{3}\left[1 - P_2(\mu)\right]$, the
potential at the point $(\mu,\ \zeta_j)$ of surface $j$ takes the
algebraic form (up to an overall sign, which cancels in the
equipotential condition)
\begin{equation}
\widetilde{U}_j(\mu, \zeta_j) \;=\;
\sum_{k=0}^{k_{\max}} P_{2k}(\mu)
\left[ A_{j,k}\, \zeta_j^{-(2k+1)} + B_{j,k}\, \zeta_j^{\,2k} \right]
\;+\; C_j\, \zeta_j^{\,2}
\;-\; \frac{1}{3}\, q\, \lambda_j^3\, \zeta_j^{\,2}
\left[ 1 - P_2(\mu) \right].
\label{eq:Utilde}
\end{equation}
The requirement that surface $j$ be an equipotential, with the equator
($\mu = 0$, $\zeta_j = 1$) chosen as the reference point, is
\citep[][Eq.~51]{Hubbard2013}
\begin{equation}
\widetilde{U}_j\big(\mu,\ \zeta_j(\mu)\big) \;=\;
\widetilde{U}_j\big(0,\ 1\big)
\qquad \text{for all } \mu \in [0,1],
\label{eq:equipot}
\end{equation}
which is one scalar nonlinear equation per $(j, \mu_a)$ pair for the
unknown $\zeta_j(\mu_a)$.  For the outermost surface ($j=0$,
$\lambda_0 = 1$) there are no exterior spheroids ($B = C = 0$), and the
monopole coefficient sums exactly to $-1$ (total mass), so the equation
simplifies to \citep[][Eq.~50]{Hubbard2013}
\begin{equation}
\frac{1}{\zeta_0} \;+\;
\sum_{k=1}^{k_{\max}} A^{(0)}_k\, P_{2k}(\mu)\, \zeta_0^{-(2k+1)}
\;+\; \frac{1}{3} q\, \zeta_0^{\,2} \left[1 - P_2(\mu)\right]
\;=\; \text{(same at $\mu=0$, $\zeta_0=1$)} .
\label{eq:outer}
\end{equation}

Two structural points deserve emphasis.  First, the physics of the
deformation lives in the \emph{monopole-with-shape} term: for interior
surfaces the $k=0$ interior term is $A_{j,0}/\zeta_j$, and
$-A_{j,0}$ is the enclosed mass fraction, so this term is
$(M_{\rm enc}/M)\,/\,\zeta_j$ --- the $1/r$ potential of the enclosed
mass evaluated on the distorted surface.  Omitting the $1/\zeta$
prefactor on the gravitational bracket (a mistake root-caused during
development; see \texttt{validation/mh24/REPORT.md}, \S7) leaves the
interior spheroids spherical and drives $J_2 \to 0$: the balance between
this term and the centrifugal $\zeta_j^2$ term \emph{is} the shape
equation.  Second, the coupling between surfaces enters only through the
coefficients $A$, $B$, $C$, which are frozen during a shape sweep; the
scheme is therefore a Jacobi-type fixed-point iteration between the
multipole quadratures (Section~\ref{sec:potential}) and the shape solves
(this section), repeated to convergence.

% ======================================================================
\section{Harmonics and Shape Observables}\label{sec:observables}
% ======================================================================

Once the shapes have converged, the conventional zonal harmonics follow
from a direct quadrature over the spheroid stack
\citep[][Eqs.~11--12]{MilitzerHubbard2024}, implemented in
\texttt{\_Jn\_from\_shapes}:
\begin{equation}
J_n \;=\; \frac{1}{M'} \sum_{i=0}^{N-1} \delta_i\, \hat{J}_{i,n},
\qquad
\hat{J}_{i,n} \;=\; -\frac{4\pi}{n+3}\, \lambda_i^{\,n+3}
\int_0^1 P_n(\mu)\, \zeta_i^{\,n+3}(\mu)\,\dd\mu ,
\label{eq:Jn}
\end{equation}
with the mass normalization
\begin{equation}
M' \;=\; \sum_i \delta_i\, \frac{4\pi}{3}\, \lambda_i^3
\int_0^1 \zeta_i^{\,3}(\mu)\,\dd\mu
\label{eq:Mnorm}
\end{equation}
(the $4\pi$ cancels in the ratio and is dropped in code).  The $n=0$
member evaluates identically to $J_0 = -1$, which serves as a built-in
sanity check on the quadrature and the mass normalization.  The
harmonics are referenced to the model's outer \emph{equatorial} radius
$a_0$, the natural CMS convention --- a point that matters when comparing
to published values referenced to a different radius
(Section~\ref{sec:saturn_rescale}).

The shape observables come directly from the converged outer surface and
the stack:
\begin{align}
\Req &= a_0, \qquad
\Rpol = a_0\, \zeta_0(\mu \to 1), \qquad
f = 1 - \zeta_0(\mu \to 1),
\label{eq:shape_obs}\\[4pt]
C_{\rm MoI} \cdot M \Req^2 \;\longleftrightarrow\;
C &= \sum_i \delta_i\, \frac{4\pi}{5}\, a_0^5\, \lambda_i^5
\int_0^1 \left(1-\mu^2\right) \zeta_i^{\,5}(\mu)\,\dd\mu ,
\label{eq:moi}
\end{align}
where $f$ is the oblateness and $C$ is the \emph{axial} (polar) moment
of inertia, obtained by applying the same density-jump superposition to
$\int \rho\, r^2 (1-\mu^2)\, \dd V$.  In the implementation the polar
value is read at the Gauss--Legendre node closest to the pole
($\mu \simeq 0.999$ for the default $L = 48$), a controlled
approximation.  The public wrappers return the ToF-compatible 8-element
vector
\begin{equation}
\texttt{j2n} \;=\;
\left[\, J_2 \times 10^6,\ J_4 \times 10^6,\ J_6 \times 10^6,\
J_8 \times 10^6,\ f,\ \Req,\ \Rpol,\ C \,\right],
\label{eq:j2n}
\end{equation}
so downstream consumers (the evolution loop, the output writer, the
scoring scripts) need not distinguish the gravity backend.

% ======================================================================
\section{Numerical Implementation}\label{sec:numerics}
% ======================================================================

\subsection{Discretization and quadrature}

The colatitude integrals in
Equations~(\ref{eq:Jint})--(\ref{eq:D}), (\ref{eq:Jn}),
and (\ref{eq:moi}) are evaluated by $L$-point Gauss--Legendre quadrature
mapped to $\mu \in [0,1]$ ($L = 48$ by default, configurable as
\texttt{cms\_nmu}).  The multipole expansion retains even orders through
$P_{2k_{\max}}$ with $k_{\max} = 10$ by default (\texttt{cms\_kmax}),
i.e., harmonics through $J_{20}$ participate in the potential even
though only $J_2$--$J_8$ are reported.  The radial discretization is set
by the caller: each input radius--density pair becomes one spheroid, so
an ORCHARD model with $N$ zones is represented at its native resolution.

\subsection{The outer iteration}

Algorithm (implemented in \texttt{cms()}):
\begin{enumerate}
\item Initialize $\zeta_i(\mu_a) = 1$ (spherical cold start), or from a
  supplied warm-start field (Section~\ref{sec:warmstart}).
\item \textbf{Multipoles:} evaluate
  Equations~(\ref{eq:Jint})--(\ref{eq:D}) from the current shapes.
\item \textbf{Coefficients:} form $A_{j,k}$, $B_{j,k}$, $C_j$,
  $A^{(0)}_k$ by cumulative sums, Equations~(\ref{eq:A})--(\ref{eq:C}).
\item \textbf{Shapes:} solve Equations~(\ref{eq:equipot}) and
  (\ref{eq:outer}) for all $(j,\mu_a)$.
\item Repeat from step 2 until
  $\max_{j,a} |\Delta\zeta| < 10^{-10}$ (at least two passes; cap 300).
\item Evaluate the harmonics and observables,
  Equations~(\ref{eq:Jn})--(\ref{eq:j2n}), from the converged shapes.
\end{enumerate}
A cold start converges in $\sim$28 Jacobi passes for a Jupiter model; a
warm start typically needs 2.

\subsection{Vectorized Newton shape solve}

The original implementation solved each of the $N \times L$ scalar
equations with \texttt{scipy.brentq}, which dominated the runtime.  The
production solver (\texttt{\_solve\_shapes}) instead runs a damped
Newton iteration with the \emph{analytic} derivative of
Equation~(\ref{eq:Utilde}) with respect to $\zeta$,
\begin{equation}
\frac{\partial \widetilde{U}_j}{\partial \zeta}
= \sum_k P_{2k}(\mu)\left[ -(2k{+}1) A_{j,k} \zeta^{-(2k+2)}
+ 2k\, B_{j,k} \zeta^{2k-1}\right]
+ 2\left( C_j - \tfrac{1}{3} q \lambda_j^3 \left[1{-}P_2\right]\right)
\zeta ,
\end{equation}
applied simultaneously to the whole $(N,L)$ shape grid.  The residual is
smooth and monotonic near the root, so Newton converges in a few passes;
iterates are clipped to $\zeta \in [0.20,\ 1.0000001]$ as a safeguard
(the upper bound admits the equator's exact $\zeta = 1$, the lower bound
excludes unphysical collapse).  Two \texttt{numba}
\texttt{@njit(parallel=True)} kernels provide the fast path
(\texttt{\_newton\_interior\_njit}, \texttt{\_newton\_outer\_njit},
and \texttt{\_multipoles\_njit} for the quadratures), with
\texttt{prange} parallelism over layers and pure-\texttt{numpy}
fallbacks retained when \texttt{numba} is unavailable.  The
vectorization contributed $\sim$65$\times$ and the JIT another
$\sim$7$\times$; together they turned a $\sim$2.3~s warm solve into
$\sim$5~ms at $N=512$ with the answers unchanged
(\texttt{validation/mh24/REPORT.md}, \S10).

\subsection{Warm starting}\label{sec:warmstart}

Both the shape field $\zeta$ (an $(N,L)$ array) and, in the wrapper of
Section~\ref{sec:meanradius}, the equatorial-radius grid are reusable
across calls.  In ORCHARD the previous call's converged $\zeta$ is held
in a module-level cache (\texttt{\_cms\_shape\_cache} in
\texttt{hydrostatic.py}), reset at \texttt{init=True} of a fresh
evolution run, mirroring the ToF4/ToF7 warm-start caches.  Because
successive hydrostatic-equilibrium calls differ only slightly, warm
starting cuts the Jacobi iteration count from $\sim$28 to $\sim$2 and is
the single largest practical speedup.

\subsection{Performance}

Like-for-like benchmarks on the same Jupiter structure
(\texttt{validation/mh24/benchmark\_gravity.py}) give, per warm-started
gravity solve: 4.7~ms (CMS) versus 2.5~ms (ToF7) at $N=512$; 8.3 versus
4.9~ms at $N=1024$; and 15.0 versus 10.0~ms at $N=2049$ --- CMS is a
factor 1.5--1.9 of ToF7, with cold starts roughly at ToF7-cold cost.
The integrated ORCHARD call (which wraps the mean-to-equatorial radius
loop of Section~\ref{sec:meanradius} around several CMS solves) costs
$\sim$180~ms at $N=512$.  CMS is therefore fully affordable for
validation and post-processing at native resolution and usable in-loop,
while ToF7 remains the cheapest default for the Gyr-scale inner loop.

% ======================================================================
\section{Integration into ORCHARD}\label{sec:integration}
% ======================================================================

\subsection{Mean-radius to equatorial-radius mapping}
\label{sec:meanradius}

ORCHARD's hydrostatic structure lives on \emph{volumetric} (mean) radii
$r_{\rm mean}$, whereas CMS is parameterized by \emph{equatorial} radii
$\lambda$.  The two are related through the converged shapes: the
volume enclosed by surface $i$ is
$\tfrac{4}{3}\pi a_0^3 \lambda_i^3 \int_0^1 \zeta_i^3\,\dd\mu$, so the
volumetric radius of the spheroid is
\begin{equation}
r_{S,i} \;=\; a_0\,\lambda_i
\left( \int_0^1 \zeta_i^{\,3}(\mu)\,\dd\mu \right)^{1/3}
\;<\; a_0 \lambda_i .
\label{eq:rS}
\end{equation}
The wrapper \texttt{get\_moments\_cms\_from\_mean} (the entry point used
by ORCHARD, matching ToF's input convention) iterates the
mass-conserving fixed point
\begin{equation}
\lambda_i \;\longleftarrow\; \lambda_i \,
\frac{r_{{\rm mean},i}}{r_{S,i}},
\label{eq:lamloop}
\end{equation}
re-solving CMS at each pass (warm-started), until the relative change is
below $10^{-7}$ (at most 6 passes; a few suffice in practice).  On the
\citet{MilitzerHubbard2024} reference model this loop recovers the
equatorial radius to $\sim$0.4~km out of 71\,492~km.

\subsection{The gravity-backend dispatcher}

The configuration key \texttt{[hydrostatic\_equilibrium]
gravity\_method = tof | cms} (default \texttt{tof}) selects the backend
inside \texttt{hydrostatic.py:\_call\_tof}, which already dispatched
\texttt{tof\_order} 4 versus 7.  All three backends return the same
8-element \texttt{j2n} of Equation~(\ref{eq:j2n}); CMS additionally
returns zero-filled placeholders for the ToF figure functions $s_{2n}$
(which it does not produce), so no call site in \texttt{evolution.py} or
\texttt{hydrostatic.py} changes.  The dispatcher passes the current
zone radii, densities, total mass, and $\omega$; reverses the arrays to
CMS's surface-to-center convention; and manages the warm-start cache.

\subsection{Rotation feedback}

Within the evolution loop, gravity solves happen inside
\texttt{get\_omega}: ORCHARD anchors the planet's angular momentum to
the present-day observed state via
\begin{equation}
\omega \;=\; \omega_{\rm obs}\,
\frac{I_{\rm planet}}{I_{\rm structure}},
\qquad I_{\rm planet} = C_{\rm MoI}\, M \Req^2 ,
\label{eq:spin}
\end{equation}
with $C_{\rm MoI}$ the observed normalized moment of inertia (0.26393
for Jupiter, 0.2181 for Saturn) and $I_{\rm structure}$ the model's
axial moment of inertia --- Equation~(\ref{eq:moi}) when CMS is the
backend.  As the planet contracts, $I_{\rm structure}$ falls and the
model spins up, feeding back into the centrifugal term of the next
hydrostatic solve.  The per-timestep harmonics and figure observables
are appended to \texttt{Js\_tof.txt} in the model output directory and
exposed by \texttt{load\_models} as \texttt{J2\_TOF}, \texttt{J4\_TOF},
\texttt{J6\_TOF}, \texttt{oblat\_TOF}, \texttt{Req\_TOF}, and
\texttt{Rpol\_TOF} time series.

% ======================================================================
\section{Validation}\label{sec:validation}
% ======================================================================

\subsection{Analytic Maclaurin oracle}

The single-spheroid limit has an exact classical solution
\citep{Chandrasekhar1969} that exercises every piece of the machinery.
A homogeneous spheroid of eccentricity $e$ has the shape and harmonics
\begin{equation}
\zeta(\mu) = \left[\,1 + \frac{e^2 \mu^2}{1 - e^2}\,\right]^{-1/2},
\qquad
J_{2n} = (-1)^{n+1}\, \frac{3\, e^{2n}}{(2n+1)(2n+3)} ,
\label{eq:maclaurin}
\end{equation}
so $J_2 = e^2/5$, $J_4 = -3e^4/35$, etc., and the equilibrium rotation
rate satisfies
\begin{equation}
\frac{\omega^2}{\pi G \rho} \;=\;
\frac{2\sqrt{1-e^2}}{e^3}\left(3 - 2e^2\right)\arcsin e
\;-\; \frac{6\left(1-e^2\right)}{e^2},
\qquad
q = \frac{3}{4}\,\frac{1}{\sqrt{1-e^2}}\,\frac{\omega^2}{\pi G \rho}.
\label{eq:chandra}
\end{equation}
The module's self-tests (\texttt{python utils/cms\_hubbard.py}) verify
(i) the $J_n$ quadrature of Equation~(\ref{eq:Jn}) on the exact shape,
which matches Equation~(\ref{eq:maclaurin}) to $\sim\!10^{-13}$, and
(ii) the full shape solver driven at the Chandrasekhar $q(e)$, which
recovers the exact ellipsoid to $\sim\!10^{-10}$ in $\zeta$ and $J_2$ to
$\sim\!10^{-8}$.

\subsection{Reproduction of Militzer \& Hubbard (2024)}

Against the published Jupiter models of \citet{MilitzerHubbard2024}
(machine-readable spheroids and harmonics from the Zenodo archive), the
solver run at their native 2049 spheroids reproduces their static
interior harmonics $J^{\rm Int}$ to relative $1.5\times10^{-5}$ ($J_2$),
$3.0\times10^{-5}$ ($J_4$), $4.5\times10^{-5}$ ($J_6$), and
$5.9\times10^{-5}$ ($J_8$), with oblateness 0.06481 against their
0.06487.  The residual shrinks as $\sim\!1/N$ with layer count and is
therefore discretization-limited, not algorithmic --- the expected
signature of representing a smooth density profile by finitely many
constant-density shells \citep[cf.\ the precision study of]
[]{DebrasChabrier2018}.  Together with ToF7 run on the same density
profiles (which agrees with CMS at the $10^{-5}$--$10^{-4}$ level,
the genuine perturbative truncation), the two independent solvers
bracket the published values and cross-validate ORCHARD's gravity
stack end to end.

% ======================================================================
\section{Application to the Gravitational Observables of Jupiter and
Saturn}\label{sec:observables_js}
% ======================================================================

CMS enters the Jupiter and Saturn programs in two modes: as the optional
in-loop gravity backend during evolution
(Section~\ref{sec:integration}), and --- the workhorse mode --- as the
\emph{CMS shell-consistent score} applied in post-processing to
evolution-grid endpoints (the \texttt{rescore\_cms\_*.py} family and
\texttt{score\_cef.py}).

\subsection{The CMS shell-consistent score}\label{sec:score}

For each evolved model, the final saved structure is read from the HDF5
output: zone-boundary radii $r_b$, zone densities $\rho$, total mass
$M$, and the final $\omega$.  Feeding the boundaries with their zone
densities to \texttt{get\_moments\_cms\_from\_mean} makes each model
shell literally one CMS spheroid --- the ``shell-consistent'' part: no
re-interpolation, the gravity field of exactly the structure the
evolution produced.  The scorer then iterates the spin to
self-consistency with the observed angular-momentum anchor,
\begin{equation}
\omega^{(n+1)} \;=\; \omega_{\rm obs}\,
\frac{C_{\rm MoI}\, M \Req^2}{C^{(n)}_{\rm CMS}} ,
\label{eq:spinloop}
\end{equation}
where $C^{(n)}_{\rm CMS}$ is the CMS axial moment of inertia of
Equation~(\ref{eq:moi}) at the current spin, warm-starting each CMS
call from the previous $\zeta$; this converges to
$2\times10^{-5}$ in a few passes.  Unlike the evolution-time
initialization guard, the score applies no clamp to the inertia ratio
(the ``honest no-clamp spin''): a model that has contracted inside the
observed moment of inertia is allowed to spin faster than
$\omega_{\rm obs}$, and pays for it in its harmonics.  A ToF4
recomputation of $J_2$ is also compared against the value saved during
the run as a consistency gate on the stored structure.  The result is a
$\chi^2$ over the observables; for Jupiter
(e.g., \texttt{score\_cef.py}) the targets are
$\Req = 71\,492 \pm 200$~km,
$J_2 \times 10^6 = 14\,696.6 \pm 30$,
$J_4 \times 10^6 = -586.6 \pm 5$ \citep{Iess2018},
$T_{\rm eff} = 125.57 \pm 1.0$~K \citep{Li2012}, and
$Y_{\rm atm} = 0.238 \pm 0.005$, with
$C_{\rm MoI} = 0.26393$ and $\omega_{\rm obs} = 2\pi/35\,730$~s
(the adopted $\sigma$'s are model-selection tolerances, deliberately far
looser than the measurement errors).

\subsection{Saturn and the gravity reference radius}
\label{sec:saturn_rescale}

The Saturn scorers (\texttt{rescore\_cms\_sat*.py}) use the same recipe
with Saturn constants ($C_{\rm MoI} = 0.2181$,
$\omega_{\rm obs} = 2\pi/38\,014$~s,
$\Req = 60\,268 \pm 170$~km,
$J_2 \times 10^6 = 16\,290.573 \pm 30$,
$J_4 \times 10^6 = -935.314 \pm 5$ from the Cassini Grand Finale
\citep{Iess2019}, $T_{\rm eff} = 96.67 \pm 1.0$~K \citep{Li2015}, and a
flat-bottomed $Y_{\rm atm}$ penalty over the Cassini UVIS+CIRS band
0.16--0.22 of \citealt{KoskinenGuerlet2018}), plus one load-bearing
addition.  Because $J_{2n}$ is defined jointly with a reference radius
(Equation~\ref{eq:external}), a harmonic quoted at radius $a$ transforms
to radius $a'$ as
\begin{equation}
J_{2n}(a') \;=\; J_{2n}(a) \left(\frac{a}{a'}\right)^{2n} .
\label{eq:rescale}
\end{equation}
CMS references the model's harmonics to its \emph{own} equatorial
radius, but the \citet{Iess2019} Saturn values are referenced to
$a' = 60\,330$~km --- not the 1-bar equatorial radius of 60\,268~km.
The Saturn score therefore rescales every model's $J_{2n}$ to the Iess
reference radius via Equation~(\ref{eq:rescale}) before comparing;
skipping this step mis-scores $J_2$ at the $\sim$5$\sigma$ level of the
adopted tolerance.  Jupiter never needed the rescale only because
Juno's reference radius coincides with the 71\,492~km target radius.

% ======================================================================
\section{Limitations}\label{sec:limitations}
% ======================================================================

\begin{enumerate}
\item \textbf{Rigid rotation.}  The centrifugal term in
  Equation~(\ref{eq:Utilde}) assumes uniform $\omega$.  Differential
  rotation and deep winds, which contribute measurably to the observed
  harmonics (the $J^{\rm Wind}$ of \citealt{MilitzerHubbard2024} and the
  wind solutions of \citealt{Iess2019}), are not modeled; ORCHARD's CMS
  output corresponds to the \emph{static interior} contribution
  $J^{\rm Int}$ and must be compared accordingly.
\item \textbf{Axisymmetry and north--south symmetry.}  Only even zonal
  harmonics are produced; odd harmonics ($J_3$, $J_5$, \ldots), which
  encode hemispheric wind asymmetries, are outside the method.
\item \textbf{Layer discretization.}  The $\sim\!1/N$ shell-discretization
  error means ORCHARD's typical $N = 500$--700 zones deliver $J_2$ to
  a few $\times 10^{-5}$ relative, versus $1.5\times10^{-5}$ at the
  2049 spheroids of \citet{MilitzerHubbard2024}.  For method-limited
  comparisons, run the final structure at higher $N$.
\item \textbf{Polar radius at a quadrature node.}  $\Rpol$ and the
  oblateness are read at the Gauss node nearest the pole
  ($\mu \simeq 0.999$), a small but nonzero approximation.
\item \textbf{Gravity only.}  CMS consumes the density--radius stack and
  $\omega$; it does not touch the EOS or thermal structure.  In the
  \citet{MilitzerHubbard2024} reproduction, the residual
  $\sim$0.8\% end-to-end $J_2$ offset traced to the H--He EOS, not to
  the gravity solver --- a reminder that at Juno/Cassini precision the
  EOS, not the figure theory, is the accuracy floor.
\end{enumerate}

% ======================================================================
\begin{thebibliography}{}

\bibitem[Chandrasekhar(1969)]{Chandrasekhar1969}
Chandrasekhar, S. 1969, Ellipsoidal Figures of Equilibrium
(New Haven, CT: Yale Univ.\ Press)

\bibitem[Debras \& Chabrier(2018)]{DebrasChabrier2018}
Debras, F., \& Chabrier, G. 2018, A\&A, 609, A97

\bibitem[Hubbard(2013)]{Hubbard2013}
Hubbard, W.~B. 2013, ApJ, 768, 43

\bibitem[Iess et al.(2018)]{Iess2018}
Iess, L., Folkner, W.~M., Durante, D., et al. 2018, Nature, 555, 220

\bibitem[Iess et al.(2019)]{Iess2019}
Iess, L., Militzer, B., Kaspi, Y., et al. 2019, Science, 364, aat2965

\bibitem[Koskinen \& Guerlet(2018)]{KoskinenGuerlet2018}
Koskinen, T.~T., \& Guerlet, S. 2018, Icarus, 307, 161

\bibitem[Li et al.(2012)]{Li2012}
Li, L., Baines, K.~H., Smith, M.~A., et al. 2012,
J.\ Geophys.\ Res.\ Planets, 117, E11002

\bibitem[Li et al.(2015)]{Li2015}
Li, L., Jiang, X., West, R.~A., et al. 2015, Geophys.\ Res.\ Lett.,
42, 2144

\bibitem[Militzer \& Hubbard(2024)]{MilitzerHubbard2024}
Militzer, B., \& Hubbard, W.~B. 2024, Icarus, 411, 115955

\bibitem[Militzer et al.(2019)]{Militzer2019}
Militzer, B., Wahl, S., \& Hubbard, W.~B. 2019, ApJ, 879, 78

\bibitem[Nettelmann(2017)]{Nettelmann2017}
Nettelmann, N. 2017, A\&A, 606, A139

\bibitem[Nettelmann et al.(2021)]{Nettelmann2021}
Nettelmann, N., Movshovitz, N., Ni, D., et al. 2021, PSJ, 2, 241

\bibitem[Wahl et al.(2017)]{Wahl2017}
Wahl, S.~M., Hubbard, W.~B., Militzer, B., et al. 2017,
Geophys.\ Res.\ Lett., 44, 4649

\bibitem[Zharkov \& Trubitsyn(1978)]{ZharkovTrubitsyn1978}
Zharkov, V.~N., \& Trubitsyn, V.~P. 1978, Physics of Planetary
Interiors (Tucson, AZ: Pachart)

\end{thebibliography}

\end{document}
